\documentclass[11pt,a4paper]{article}
\usepackage[utf8]{inputenc}
\usepackage[T1]{fontenc}
\usepackage[portuguese]{babel}
\usepackage{geometry}
\geometry{margin=2.5cm}
\usepackage{listings}
\usepackage{xcolor}
\usepackage{amsmath}
\usepackage{amssymb}
\usepackage{hyperref}
\usepackage{enumitem}
\usepackage{tcolorbox}
\usepackage{tabularx}
\usepackage{booktabs}
\usepackage{graphicx}

\definecolor{codebg}{rgb}{0.95,0.95,0.95}
\definecolor{codegreen}{rgb}{0.1,0.5,0.1}
\definecolor{codepurple}{rgb}{0.5,0.1,0.5}
\definecolor{codegray}{rgb}{0.5,0.5,0.5}

\lstdefinestyle{python}{
    backgroundcolor=\color{codebg},
    commentstyle=\color{codegray}\itshape,
    keywordstyle=\color{codepurple}\bfseries,
    stringstyle=\color{codegreen},
    basicstyle=\ttfamily\small,
    breaklines=true,
    frame=single,
    language=Python,
    showstringspaces=false,
    tabsize=4,
    literate={á}{{\'a}}1 {ã}{{\~a}}1 {é}{{\'e}}1 {í}{{\'i}}1 {ó}{{\'o}}1 {õ}{{\~o}}1 {ú}{{\'u}}1 {ç}{{\c{c}}}1 {Á}{{\'A}}1 {Ã}{{\~A}}1 {É}{{\'E}}1 {Í}{{\'I}}1 {Ó}{{\'O}}1 {Õ}{{\~O}}1 {Ú}{{\'U}}1 {Ç}{{\c{C}}}1 {â}{{\^a}}1 {ê}{{\^e}}1 {ô}{{\^o}}1 {û}{{\^u}}1 {Â}{{\^A}}1 {Ê}{{\^E}}1 {Ô}{{\^O}}1 {Û}{{\^U}}1 {à}{{\`a}}1 {è}{{\`e}}1 {ì}{{\`i}}1 {ò}{{\`o}}1 {ù}{{\`u}}1 {ä}{{\"a}}1 {ö}{{\"o}}1 {Ä}{{\"A}}1 {Ö}{{\"O}}1
}
\lstset{style=python}

\newtcolorbox{objetivos}[1][]{
  colback=green!5!white,
  colframe=green!40!black,
  title=O que você vai aprender nesta página,
  fonttitle=\bfseries,
  #1
}

\newtcolorbox{exercicio}[1]{
  colback=blue!5!white,
  colframe=blue!40!black,
  title=#1,
  fonttitle=\bfseries
}

\newtcolorbox{solucao}{
  colback=yellow!5!white,
  colframe=orange!60!black,
  title=Solução proposta,
  fonttitle=\bfseries
}

\title{\textbf{Data Analysis with Python 2024--2025}\\Parte 6: Machine Learning Types\\\large Caderno de Exercícios}
\author{Tradução e Adaptação: University of Helsinki --- MOOC.fi}
\date{}

\begin{document}

\maketitle

\section*{Nota Importante}
Este material é uma tradução e adaptação baseada no curso \textit{Data Analysis with Python 2024-2025}, oferecido pela \textbf{Universidade de Helsinque (MOOC.fi)}.

\tableofcontents
\newpage

\section{Exercícios - Capítulo 6 (Clustering)}

\subsection*{Exercício 6.5 -- Clustering de Plantas (plant\_clustering)}

\begin{exercicio}{Sumário}
\begin{itemize}
    \item \textbf{Objetivo:} Aplicar o agrupamento K-Means no clássico conjunto de dados Iris com 3 clusters definidos.
    \item \textbf{Procedimento:}
    \begin{enumerate}
        \item Crie a função \texttt{plant\_clustering()} que carregue os dados (\texttt{load\_iris}).
        \item Inicialize o classificador \texttt{KMeans(3, random\_state=0)} e faça o ajuste aos dados de treinamento.
        \item Use a técnica de permutação de rótulos (\texttt{find\_permutation}) introduzida na lição para realinhar os nomes arbitrários dos clusters aos números corretos baseando-se na moda do alvo correto.
    \end{enumerate}
    \item \textbf{Retorno:} A \texttt{accuracy\_score} da correspondência entre os rótulos do k-means com permutações e o alvo biológico real.
\end{itemize}
\end{exercicio}

\subsection*{Exercício 6.6 -- Clusters Não Convexos (nonconvex\_clusters)}

\begin{exercicio}{Sumário}
\begin{itemize}
    \item \textbf{Objetivo:} Ler o arquivo \texttt{src/data.tsv}, processá-lo pelo \texttt{DBSCAN} variando o limiar de busca (\texttt{eps}) e exportar uma matriz de relatórios do modelo em formato pandas \texttt{DataFrame}.
    \item \textbf{Procedimento:}
    \begin{enumerate}
        \item Realize um laço para executar instâncias do \texttt{DBSCAN} baseados nos hiperparâmetros gerados por \texttt{np.arange(0.05, 0.2, 0.05)}.
        \item \textbf{Importante:} O DBSCAN rotula anomalias/outliers isoladas como \texttt{-1}. Você deve adaptar a função \texttt{find\_permutation} para não computá-los durante a contagem de acurácia.
        \item Caso o modelo crie um número final de agrupamentos não corrompidos que destoe da quantidade de rótulos originais dos dados (após excluir os outliers), marque a acurácia para aquela etapa como \texttt{np.nan}.
    \end{enumerate}
    \item \textbf{Retorno:} Um \texttt{DataFrame} cujas colunas devem ser estritamente: \texttt{["eps", "Score", "Clusters", "Outliers"]}.
\end{itemize}
\end{exercicio}

\subsection*{Exercício 6.7 -- Locais de Ligação Molecular (binding\_sites)}

\begin{exercicio}{Sumário}
\begin{itemize}
    \item \textbf{Objetivo:} Construir um fluxo usando o \texttt{AgglomerativeClustering} (agrupamento hierárquico) capaz de discernir conjuntos genéticos (Adenina, Citosina, Guanina e Timina) a partir do \texttt{data.seq}.
    \item \textbf{Parte 1 (\texttt{toint} e \texttt{get\_features\_and\_labels}):} Escreva a função que converte o nucleotídeo (A, C, G, T) num escalar sequencial (0, 1, 2, 3) iterando sobre a string contida na coluna X do dataset para criar a sua matriz de \textit{features}.
    \item \textbf{Parte 2 (\texttt{cluster\_euclidean}):} Crie um agrupamento aglomerativo hierárquico, estipulando 2 clusters, métrica de afinidade \textit{euclidiana} com critério de união \emph{average} (ligação por média). Use a permutação para recuperar rótulos e retorne o \emph{score}.
    \item \textbf{Parte 3 (\texttt{cluster\_hamming}):} Substitua a forma base e agora aplique a contagem de instâncias pela \textit{distância de Hamming}. Utilize a função explícita \texttt{pairwise\_distances(..., metric='hamming')} sob a matriz principal, e crie a classe de agrupamento instruída agora para aceitar vetores pré-calculados (\texttt{affinity='precomputed'}). Retorne a acurácia.
\end{itemize}
\end{exercicio}

\end{document}
